\section{CNN Hyperparameter Search}
\label{sec:cnn-hyperparameters}

\subsection{Search Space}

The tuning grid is stored in \texttt{configs/cnn.yaml} and read by \texttt{python -m signbridge.tune\_cnn}. Five hyperparameters are varied: the head learning rate, the batch size, the dropout probability, the AdamW weight decay and the number of unfrozen feature modules. The grid in Table~\ref{tab:cnn-search-space} contains $2^5=32$ combinations. Training is capped at 12 trials, and each trial runs for 3 epochs. The 12 trials are a seeded random sample, not the first 12 rows of the grid: \texttt{random.Random} with the config seed (42) shuffles all 32 combinations, duplicate combinations are removed, and the search walks that shuffled order from the start and runs at most 12 trials. A position whose parameter set and seed are already in the trials file is skipped, so an interrupted search fills the remaining budget in the same order instead of starting a new search. Early stopping is configured with patience 3, so it does not trigger inside a 3-epoch trial. These short trials are intended to rank configurations, not to produce the final model.

\begin{table}[htbp]
\centering
\caption{CNN hyperparameter search space in \texttt{configs/cnn.yaml}.}
\label{tab:cnn-search-space}
\begin{tabular}{lll}
\toprule
Hyperparameter & Values & Role \\
\midrule
Head learning rate & $3\times10^{-4}$, $1\times10^{-3}$ & step size for the new classifier \\
Batch size & 32, 64 & images per optimisation step \\
Dropout & 0.2, 0.4 & regularisation before the linear layer \\
Weight decay & $1\times10^{-4}$, $1\times10^{-3}$ & AdamW regularisation \\
Unfrozen blocks & 0, 2 & last feature modules trained \\
\midrule
Trials & at most 12 of 32 & seeded random sample, hard cap \\
Epochs per trial & 3 & short ranking runs \\
\bottomrule
\end{tabular}
\end{table}

The remaining settings are fixed during tuning: 26 classes, ImageNet-pretrained MobileNetV2, $224\times224$ input, cross-entropy with label smoothing 0.1, AdamW with a backbone learning rate multiplier of 0.1, cosine annealing, early-stopping patience 3 and seed 42. The final run uses the same configuration with 12 epochs and early stopping; the real run stopped after epoch 9.

\subsection{Selection Rule}

A trial is ranked by its best validation macro-F1. Ties are broken in favour of the trial with fewer trainable parameters, because a smaller model is cheaper to train and to run and is less likely to overfit. The rule is written in the docstring of \texttt{tune\_cnn.py} and in \texttt{DECISIONS.md}. The test split is never read during tuning. Each trial appends one row to \texttt{reports/tuning/cnn\_trials.csv}; a parameter set and seed that already appear in that file are skipped, so an interrupted search can be resumed without repeating work.

\subsection{Trial Results}

\IfFileExists{../../reports/tuning/cnn_trials_table.tex}{%
  \begin{table}[htbp]
  \centering
  \caption{Tuning trials for the CNN, ordered as recorded. The table is generated from \texttt{reports/tuning/cnn\_trials.csv} by \texttt{tools/trials\_to\_tex.py}.}
  \label{tab:cnn-trials}
  \resizebox{\linewidth}{!}{\begin{tabular}{lrrrrrrrrr}
\hline
run\_name & lr & batch\_size & dropout & weight\_decay & unfreeze\_blocks & best\_val\_macro\_f1 & val\_accuracy & epochs\_run & seconds \\
\hline
trial\_000 & 0.001 & 64 & 0.2 & 0.001 & 0 & 0.8065300626970106 & 0.8108974358974359 & 3 & 132.15 \\
trial\_001 & 0.0003 & 32 & 0.4 & 0.0001 & 2 & 0.8794509384109741 & 0.8801282051282051 & 3 & 125.61 \\
trial\_002 & 0.0003 & 64 & 0.2 & 0.001 & 0 & 0.6753639935073358 & 0.6820512820512821 & 3 & 129.19 \\
trial\_003 & 0.0003 & 64 & 0.4 & 0.001 & 2 & 0.8430798718360818 & 0.8442307692307692 & 3 & 132.34 \\
trial\_004 & 0.001 & 64 & 0.2 & 0.0001 & 2 & 0.9225361978376546 & 0.9224358974358975 & 3 & 131.39 \\
trial\_005 & 0.0003 & 64 & 0.2 & 0.001 & 2 & 0.8527992548457175 & 0.8538461538461538 & 3 & 131.39 \\
trial\_006 & 0.001 & 32 & 0.4 & 0.001 & 0 & 0.8288342444338248 & 0.8314102564102565 & 3 & 126.31 \\
trial\_007 & 0.0003 & 32 & 0.4 & 0.001 & 0 & 0.7260288440272608 & 0.732051282051282 & 3 & 124.31 \\
trial\_008 & 0.001 & 32 & 0.2 & 0.001 & 2 & 0.9437027962531004 & 0.9442307692307692 & 3 & 127.6 \\
trial\_009 & 0.0003 & 64 & 0.4 & 0.0001 & 0 & 0.6644214014818777 & 0.6705128205128205 & 3 & 128.42 \\
trial\_010 & 0.001 & 32 & 0.2 & 0.0001 & 0 & 0.8364909320199915 & 0.8384615384615385 & 3 & 123.65 \\
trial\_011 & 0.0003 & 64 & 0.2 & 0.0001 & 2 & 0.8527992548457175 & 0.8538461538461538 & 3 & 132.25 \\
\hline
\end{tabular}
}
  \end{table}
}{%
  \todobox{The CNN trial table is missing. It will be generated from \texttt{reports/tuning/cnn\_trials.csv} by \texttt{tools/trials\_to\_tex.py} after the real tuning runs on the shared split.}
}

All 12 planned trials ran, each for three epochs, on the shared split. Together they took 1{,}545 seconds of training, about 26 minutes, on the cloud GPU. Table~\ref{tab:cnn-trials} lists every trial. The winner is \texttt{trial\_008}: head learning rate $1\times10^{-3}$, batch size 32, dropout 0.2, weight decay $1\times10^{-3}$ and the last two feature modules unfrozen. It reached a validation accuracy of 0.9442 and a validation macro-F1 of 0.9437 after three epochs. The runner-up, \texttt{trial\_004}, used the same learning rate and dropout with batch size 64, weight decay $1\times10^{-4}$ and two unfrozen modules, and reached a validation accuracy of 0.9224 and a macro-F1 of 0.9225. The winner was unique, so the tie-break rule did not decide it. Two lower trials, \texttt{trial\_005} and \texttt{trial\_011}, did tie on a best validation macro-F1 of 0.8528; both have the same number of trainable parameters, so the tie-break rule would not separate them either.

The spread over the 12 trials is large: the worst trial, \texttt{trial\_009}, reached 0.6644, so the range is 0.2793. The clearest pattern is the number of unfrozen modules. Every trial that trained the head only (0 unfrozen modules) scored at most 0.8365, while all six trials with two unfrozen modules scored at least 0.8431; the mean rises from 0.7563 to 0.8824. Fine-tuning the last two blocks is therefore the choice that matters most. The learning rate also separates the trials: the five runs at $1\times10^{-3}$ average 0.8676 against 0.7848 for the seven at $3\times10^{-4}$, although the grid changes several values at once, so this describes the sample rather than an isolated effect.

The differences should be read with care. Each configuration was run once, so the table contains no estimate of seed-to-seed noise. The winner's validation accuracy is 0.9442 against 0.9224 for the runner-up, a difference of 34 of the 1{,}560 validation images, but with one run per configuration this is not a significance test. The winner was chosen by the selection rule above before the test split was read. Its values were copied into \texttt{configs/cnn.yaml} for the final run, and the first epoch of that run reproduced the winner's first-epoch training loss (1.549029) and training accuracy (0.728434) exactly, which confirms that the final run used the winning settings.
